
\section{Problema, objetivos y requisitos}
\subsection{Problema y alcance}
La empresa necesita que una venta mantenga coherentes el pedido el pago y las existencias. Un kit puede contener otros kits por lo que el stock debe calcularse sobre componentes físicos. Los reintentos de cobro o notificación no pueden duplicar salidas de inventario ni reembolsos. Además los cambios posteriores del catálogo no deben reescribir ventas ya registradas.

El alcance corresponde al caso académico de la guía. No se implementan pagos bancarios reales envíos externos de correo o SMS facturación tributaria ni cálculos fiscales no definidos \cite{guia}.

\subsection{Objetivos}
\textbf{General} Diseñar implementar probar y documentar una aplicación de escritorio con persistencia y trazabilidad que integre Clean Code SOLID Clean Architecture y los 23 patrones GoF en procesos verificables

\textbf{Específicos}
\begin{enumerate}[label=O\arabic*,leftmargin=1.2cm]
\item Modelar actores requisitos invariantes y relaciones de ventas e inventario
\item Separar los cinco proyectos y comprobar la dirección de sus referencias
\item Integrar las tres familias GoF en operaciones reales y localizar sus participantes
\item Proteger confirmación devolución y eventos frente a fallos duplicados y concurrencia
\item Validar resultados con pruebas matriz de aceptación diagramas Git y evidencias reproducibles
\end{enumerate}

\subsection{Actores y casos de uso}
\begingroup\small
\noindent
\begin{tabular}{@{}>{\raggedright\arraybackslash}p{0.15\linewidth}>{\raggedright\arraybackslash}p{0.77\linewidth}@{}}
\toprule
\textbf{Actor} & \textbf{Operaciones visibles y autorizadas}\\
\midrule
ADMIN & Administración completa del catálogo clientes ventas bodega reportes auditoría y opciones de demostración\\[3pt]
VENTAS & Clientes carrito pedidos validación cobro devoluciones y reportes Catálogo de lectura\\[3pt]
BODEGA & Catálogo de lectura inventario movimientos despacho y entrega\\[3pt]
CONSULTA & Panel catálogo clientes pedidos seguimiento devoluciones inventario y reportes sin escritura de negocio\\[3pt]
\bottomrule
\end{tabular}
\par\endgroup

El menú omite apartados y acciones según el rol pero el servicio vuelve a autorizar las operaciones. No se confunde ocultar un botón con aplicar una política de acceso. La consulta y exportación de comprobantes no alteran el pedido ni el inventario.

\modelo{contexto}{Contexto del sistema elaborado a partir del código}{contexto}
\FloatBarrier
\subsection{Requisitos funcionales}
\begingroup\small
\begin{longtable}{@{}>{\raggedright\arraybackslash}p{0.09\linewidth}>{\raggedright\arraybackslash}p{0.4\linewidth}>{\raggedright\arraybackslash}p{0.43\linewidth}@{}}
\toprule
\textbf{ID} & \textbf{Requisito} & \textbf{Verificación}\\
\midrule
\endfirsthead
\toprule
\textbf{ID} & \textbf{Requisito} & \textbf{Verificación}\\
\midrule
\endhead
RF01 & Acceso local, roles, cierre y hash & Seguridad y pruebas por rol\\[3pt]
RF02 & CRUD de productos e integridad de SKU & ServicioCatalogo\\[3pt]
RF03 & Kits anidados sin ciclos & P18, P27, P34\\[3pt]
RF04 & Clientes con categoria y estado & ClienteInactivoBloqueaValidacion\\[3pt]
RF05 & Borradores, lineas y cantidades & P13, P14, P29\\[3pt]
RF06 & Descuentos y entrega & P02-P04, P12, P32\\[3pt]
RF07 & Validacion y ciclo de estados & P05, P08, P09, P23, P31\\[3pt]
RF08 & Cobro simulado e idempotencia & P06, P07, P16, P19, P22\\[3pt]
RF09 & Inventario trazable y concurrencia & P02-P07, P36\\[3pt]
RF10 & Despacho, guia y entrega & P09, P12\\[3pt]
RF11 & Devolucion total y reembolso & P33, P35\\[3pt]
RF12 & Outbox, avisos y auditoria & P11, P30, P37\\[3pt]
RF13 & Filtros y comprobantes & P17, P25, P26\\[3pt]
RF14 & Persistencia y seed repetible & P01, P10\\[3pt]
RF15 & Doce areas visuales y Mediator & Revision visual, P28\\[3pt]
\bottomrule
\end{longtable}
\endgroup

\subsection{Requisitos no funcionales medibles}
\begingroup\small
\begin{longtable}{@{}>{\raggedright\arraybackslash}p{0.1\linewidth}>{\raggedright\arraybackslash}p{0.25\linewidth}>{\raggedright\arraybackslash}p{0.57\linewidth}@{}}
\toprule
\textbf{ID} & \textbf{Criterio} & \textbf{Medida o condición}\\
\midrule
\endfirsthead
\toprule
\textbf{ID} & \textbf{Criterio} & \textbf{Medida o condición}\\
\midrule
\endhead
RNF01 & Integridad & Stock nunca negativo; un pago y devolución por pedido\\[3pt]
RNF02 & Seguridad local & Hash con sal; permisos verificados sin UI\\[3pt]
RNF03 & Mantenibilidad & Cinco proyectos y cero referencias prohibidas\\[3pt]
RNF04 & Testabilidad & Al menos 25 pruebas automaticas ejecutadas\\[3pt]
RNF05 & Recuperabilidad & Reintento Outbox sin duplicar pago ni aviso\\[3pt]
RNF06 & Trazabilidad & Fecha, usuario, resultado y correlacion en auditoria\\[3pt]
RNF07 & Usabilidad & TAB, etiquetas, foco y confirmaciones: criterio de revisión manual pendiente de cierre completo\\[3pt]
RNF08 & Rendimiento de representacion & 1000 filas Flyweight con identidad y asignaciones comparables\\[3pt]
RNF09 & Reproducibilidad & Migraciones y seed desde archivo inexistente\\[3pt]
RNF10 & Exactitud monetaria & decimal, dos decimales y AwayFromZero\\[3pt]
\bottomrule
\end{longtable}
\endgroup

No se declara un tiempo máximo de respuesta que no haya sido medido. La evidencia de 1000 filas Flyweight compara asignaciones bajo condiciones equivalentes y no memoria total del proceso. El requisito de accesibilidad conserva la revisión manual como pendiente.
